\textbf{Objective:} To describe carbon-plate insole use during preparation and competition among competitive Italian athletics athletes, estimate its cross-sectional, period-matched association with self-reported current-season injury, and examine whether that association was better characterized by relative carbon share or absolute carbon-use frequency.

\textbf{Design:} Cross-sectional survey with two-phase (voluntary, then targeted-recall) recruitment, poststratified to national age--sex population strata; reported in line with STROBE and the STROBE Sports Injury and Illness Surveillance (STROBE-SIIS) extension.

\textbf{Participants:} FIDAL-affiliated athletics athletes (Under 20, Under 23, Senior categories) recruited nationally in 2026. Of 355 eligible, complete respondents, 352 (704 athlete-period records) formed the balanced analysis cohort.

\textbf{Exposure:} Carbon share (the proportion of reported weekly training-type frequency delivered in carbon-plated spikes) and absolute carbon-frequency score (the summed carbon-frequency items) across five training domains, calculated separately for preparation and competition.

\textbf{Outcome:} Self-reported injury occurring in the preparation period and/or the competition period of the current season.

\textbf{Statistical methods:} Bayesian multilevel logistic models used period-specific intercepts, an athlete-level random intercept, and adjustment for demographic, anthropometric, period-matched training-frequency, and discipline covariates. The primary model represented carbon share among users with a continuous, data-driven curve fitted by a Hilbert-space Gaussian process, letting the amount of curvature in the dose--response relationship be estimated rather than fixed in advance. Two matched linear models then evaluated competing relative-use (carbon share) and absolute-use (carbon-frequency score) hypotheses on a common among-user standard-deviation scale. Estimands included adjusted prevalence contrasts and conditional absolute-risk differences for a fixed reference profile.

\textbf{Results:} Carbon-plated spikes accounted for a mean 35.7\% of reported preparation training and 47.7\% of reported competition training. Adoption showed no clear association with injury in either period. The primary curve showed adjusted competition-period injury prevalence climbing steadily with carbon share; moving from the 25th to 75th percentile of competition-period carbon share was associated with a 13.2-percentage-point higher adjusted injury prevalence (95\% CrI $-0.2$ to 29.1; $P_{\mathrm{dir}}=96.9\%$), the strongest of the four primary contrasts and the only one clearing the pre-specified practical-relevance threshold. In the matched linear comparison, a one-standard-deviation increase meant 24.1 percentage points of carbon share (37.3\% to 61.4\%) or 3.36 units of the summed ordinal carbon-frequency score (2.83 to 6.19). Among competition users, these increments corresponded to risk differences of 7.9 points for share (95\% CrI 1.4 to 14.4; $P_{\mathrm{dir}}=99.0\%$) and 8.3 points for absolute score (95\% CrI 1.6 to 15.9; $P_{\mathrm{dir}}=99.2\%$); preparation-period estimates were 3.6 points (95\% CrI $-3.9$ to 11.2; $P_{\mathrm{dir}}=82.8\%$) and 1.5 points (95\% CrI $-6.9$ to 9.4; $P_{\mathrm{dir}}=63.2\%$), respectively.

\textbf{Conclusion:} Among competition users, greater reported carbon exposure was associated with higher adjusted injury prevalence whether expressed as relative share or absolute frequency score; these data did not identify one scale as the better explanation. Adoption and preparation-period exposure were not clearly associated with injury. The cross-sectional, retrospective design does not support causal interpretation.
